\begin{derivation}[从\eqnrefs{eq:ee-qqg-current-dp68}到\eqnrefs{eq:ee-threejet-distribution-dp14}]
取 $K^2=s$，并定义
\begin{equation*}
s_{12}=2p_1\!\cdot p_2,\qquad
s_{1g}=2p_1\!\cdot k,\qquad
s_{2g}=2p_2\!\cdot k,\qquad
s=s_{12}+s_{1g}+s_{2g}.
\end{equation*}
无质量条件 $p_1^2=p_2^2=k^2=0$ 使两条内部夸克传播子分别为 $s_{1g}^{-1}$ 与 $s_{2g}^{-1}$。把正文的强子流写成
\begin{equation*}
J_g^{\mu\rho,a}
=g_s\,\bar u(p_1)T^a\mathcal A^{\mu\rho}v(p_2),
\end{equation*}
其中
\begin{equation*}
\mathcal A^{\mu\rho}
=\frac{\gamma^\rho(\slashed p_1+\slashed k)\gamma^\mu}{s_{1g}}
-\frac{\gamma^\mu(\slashed p_2+\slashed k)\gamma^\rho}{s_{2g}}.
\end{equation*}
先检验外胶子的 Ward 恒等式。把 $\varepsilon_\rho^*$ 换成 $k_\rho$，第一项用
\begin{equation*}
\slashed k(\slashed p_1+\slashed k)
=2k\!\cdot p_1-\slashed p_1\slashed k
\end{equation*}
并左乘 $\bar u(p_1)$，得到 $+\bar uT^a\gamma^\mu v$；第二项同理给 $-\bar uT^a\gamma^\mu v$，所以
\begin{equation*}
k_\rho J_g^{\mu\rho,a}=0.
\end{equation*}
因此胶子物理极化求和中的规范纵向项不贡献，可以用 $-g_{\rho\sigma}$；电磁流守恒同样允许对虚光子的横向部分使用 $-g_{\mu\nu}$。

把三体结果除以同一虚光子产生 $q\bar q$ 的 Born 张量，颜色部分只留下
\begin{equation*}
\frac1{N_c}\sum_{a,i,j}(T^a)_{ij}(T^a)_{ji}
=\frac1{N_c}\sum_a\Tr(T^aT^a)=C_F.
\end{equation*}
自旋部分可统一写成一个迹：
\begin{equation*}
\mathcal T_3
=-g_{\mu\nu}(-g_{\rho\sigma})
\Tr\!\left[
\slashed p_1\mathcal A^{\mu\rho}
\slashed p_2\,\overline{\mathcal A}^{\nu\sigma}
\right],
\qquad
\overline{\mathcal A}^{\nu\sigma}
\equiv\gamma^0(\mathcal A^{\nu\sigma})^\dagger\gamma^0.
\end{equation*}
展开后共有两个平方项和两个干涉项。每个偶数 $\gamma$ 迹都用递推恒等式
\begin{equation*}
\Tr(\slashed a_1\cdots\slashed a_{2n})
=\sum_{j=2}^{2n}(-1)^j\,2a_1\!\cdot a_j\,
\Tr(\slashed a_2\cdots\widehat{\slashed a_j}\cdots\slashed a_{2n})
\end{equation*}
降到 $\Tr(\slashed a\slashed b)=4a\!\cdot b$。再用
\begin{equation*}
2p_1\!\cdot p_2=s_{12},\qquad
2p_1\!\cdot k=s_{1g},\qquad
2p_2\!\cdot k=s_{2g}
\end{equation*}
消去全部标量积，平方项和干涉项合并为标准的夸克--反夸克 antenna：
\begin{equation*}
\frac{\overline{|\mathcal M_{q\bar qg}|^2}}
{\overline{|\mathcal M_{q\bar q}|^2}}
=2g_s^2C_F
\left[
\frac{s_{1g}}{s\,s_{2g}}
+\frac{s_{2g}}{s\,s_{1g}}
+\frac{2s_{12}}{s_{1g}s_{2g}}
\right].
\end{equation*}
这一行仍保留了三个物理来源：前两项分别在 $p_2\parallel k$ 与 $p_1\parallel k$ 时产生共线极点，最后一项是两条辐射路径的相干干涉。

在质心系定义 $x_i=2E_i/E_{\rm cm}$，$E_{\rm cm}=c\sqrt s$。由
\begin{equation*}
(K-p_1)^2=s_{2g}=s(1-x_1),\qquad
(K-p_2)^2=s_{1g}=s(1-x_2),
\end{equation*}
以及 $x_1+x_2+x_3=2$，有
\begin{equation*}
s_{12}=s(1-x_3)=s(x_1+x_2-1).
\end{equation*}
于是 antenna 括号逐项化为
\begin{align*}
&s\left[
\frac{s_{1g}}{s\,s_{2g}}
+\frac{s_{2g}}{s\,s_{1g}}
+\frac{2s_{12}}{s_{1g}s_{2g}}
\right]\\
&\qquad=
\frac{(1-x_2)^2+(1-x_1)^2+2(x_1+x_2-1)}{(1-x_1)(1-x_2)}
=\frac{x_1^2+x_2^2}{(1-x_1)(1-x_2)}.
\end{align*}

最后计算相空间，而不是把 Jacobian 当作已知常数。对三个无质量末态，在去掉整体空间取向后，Dalitz 变量给出
\begin{equation*}
\dd\Phi_3
=\frac{1}{128\pi^3s}\,\dd s_{1g}\,\dd s_{2g}
=\frac{s}{128\pi^3}\,\dd x_1\,\dd x_2,
\end{equation*}
物理区域由 $s_{ij}>0$ 给出
\begin{equation*}
0<x_1<1,\qquad0<x_2<1,\qquad x_1+x_2>1.
\end{equation*}
同一总不变量质量下的无质量二体相空间积分为 $\Phi_2=1/(8\pi)$，所以
\begin{equation*}
\frac{\dd\Phi_3}{\Phi_2}
=\frac{s}{16\pi^2}\,\dd x_1\,\dd x_2.
\end{equation*}
把它与上面的振幅平方比相乘，并用 $g_s^2=4\pi\alpha_s$，得到
\begin{equation*}
\frac1{\sigma_0}
\frac{\dd^2\sigma}{\dd x_1\dd x_2}
=\frac{\alpha_s}{2\pi}C_F
\frac{x_1^2+x_2^2}{(1-x_1)(1-x_2)},
\end{equation*}
即正文\eqnrefs{eq:ee-threejet-distribution-dp14}。归一化到 $\sigma_0$ 后，Born 电磁流、入射通量与共同的二体归一化因子代数相消，剩余因子正是 QCD 辐射核与三体相空间 Jacobian。
\end{derivation}
